\Askhsh{27322}{4}
\begin{ylh}{1}
    \item 1.7 Όρια συνάρτησης στο άπειρο
    \item 2.2 Παραγωγίσιμες συναρτήσεις - Παράγωγος συνάρτηση
    \item 2.3 Κανόνες παραγώγισης
    \item 3.5 Η συνάρτηση [ορισμένο ολοκλήρωμα της f από α έως χ].
\end{ylh}
ΘΕΜΑ 4

Ο νόμος του Νεύτωνα που αφορά την μείωση της θερμοκρασίας $T$ (σε βαθμούς Κελσίου) ενός σώματος συναρτήσει του χρόνου $t$ (σε ώρες), ορίζεται από την εξίσωση
\[ T(t) = E + (T_0 - E)e^{-kt} \]
όπου:
\begin{itemize}
    \item $E$ είναι η σταθερή θερμοκρασία του περιβάλλοντος χώρου στον οποίο βρίσκεται το σώμα με $E < T_0$.
    \item $T_0 = T(0)$ είναι η αρχική θερμοκρασία του σώματος τη στιγμή που τοποθετείται στο περιβάλλοντα χώρο.
    \item $k$ είναι μια θετική σταθερά.
\end{itemize}

\begin{alist}
    \item Να υπολογίσετε το $\lim_{t \to +\infty} T(t)$ και να ερμηνεύσετε το αποτέλεσμα. \monades{8}
    \item Να αποδείξετε ότι $T'(t) = k[E - T(t)]$. \monades{7}
    \item Να αποδείξετε ότι το ολοκλήρωμα $I = \int_0^1 (E - T(t)) \cdot \ln(T(t))dt$ ισούται με $\frac{2e^3-3e^4}{k}$ αν είναι $T(0) = e^4$ και $T(1) = e^3$. \monades{10}
\end{alist}